\chapter{Opcode index}\label{app:c}

\Cref{tab:opcode-index} lists every opcode this document mentions, as \href{https://github.com/sbryngelson/AGXForge/blob/main/isa/g17-opcode-glossary.json}{the opcode glossary} records it. The names are
this project's and describe the measured function; Apple's spellings are not used. \emph{Standing} here is the
glossary's own confidence scale (MM~\mm{0.16}). It can lag the newer sections that \cref{app:b} reads, and the daggers
there mark such opcodes.
\begin{itemize}
\item \emph{executed}: the behavior confirmed by running it on this GPU.
\item \emph{compiled}: seen in compiled code, including programs that ran correctly, without its own semantics
  separately confirmed.
\item \emph{static}: known only from Apple's tables or the corpus.
\item \emph{unconfirmed}: a partial identification whose name says what is known and marks what is not.
\end{itemize}
The two vocabularies are different scales, not different spellings. The closest counterparts: \cref{app:b}'s
\emph{measured} is \emph{executed} here (its \emph{measured (receipt fit)} is weaker, a replay that is not independent
evidence); its \emph{whole-program} and \emph{Apple's compiler emits} are \emph{compiled}. \emph{static} and \emph{unconfirmed} have no counterpart there, because \cref{app:b} lists only
opcodes whose function the sources state.
A few opcodes the chapters mention have no glossary entry; the table marks them, and the chapter that mentions one says what is known.
"Discussed in" gives the one or two sections that treat the opcode most; the last column gives up to four machine-model sections.

% The table is generated on every build from the glossary and from where the chapters mention each opcode
% (tools/g17docs.py --write-opindex, run by docs/tex/Makefile).
\input{opindex}
